% ---------------------------------------------------

\section{Related Work}
\label{sec:relatedWork}

Autonomous sailing has been an active research area for over a decade. Stekzer's dissertation on autonomous sailboat navigation set foundational algorithms for waypoint tracking and tacking algorithms under varying wind conditions. From that research, we found that a state-machine approach to points of sail is most efficient in open water [1]. 

Commercial companies such as Saildrone have validated the viability of autonomous wind-driven vessels for long-endurance ocean monitoring, using rigid wingsail designs similar to ALAN's [2]. 

Furthermore, Sun et al. provides a comprehensive review of autonomous sailboat sail designs, concluding that rigid wingsails consistently outperform cloth sails in autonomous applications due to their repeatability, durability, and ability to be precisely controlled with low-power actuators [7]. 

Also, Skene's foundational work on yacht design informs ALAN's hull geometry decisions, particularly the relationship between center of effort and center of lateral resistance [6]. 

Current autonomaus platforms are either expensive (Saildrone) or limited to single-hull monohull designs with conventional cloth sails. ALAN addresses this gap by combining a trimaran hull for superior stability with a self-trimming rigid wingsail, built entirely from accessible components and low-cost fabrication resources.

% ---------------------------------------------------
